% =====================================================================
%  preamble.tex -- shared template for the two placement study books
%
%  Large parts of this file are adapted from the preamble and macros of
%  Evan Chen's "An Infinitely Large Napkin"
%  (https://github.com/vEnhance/napkin, files tex/preamble.tex and
%  tex/macros.tex).  Those source files are licensed under the
%  GNU General Public License, version 3.  Accordingly this adapted file
%  is also distributed under the GNU GPL v3; see latex/COPYING-GPL-3.txt
%  and latex/NOTICE.md.
%
%  Copyright (C) Evan Chen and the Napkin contributors (original code).
%  Modifications (C) Philospher.
%
%  This program is free software: you can redistribute it and/or modify
%  it under the terms of the GNU General Public License as published by
%  the Free Software Foundation, either version 3 of the License, or
%  (at your option) any later version.  It is distributed WITHOUT ANY
%  WARRANTY; without even the implied warranty of MERCHANTABILITY or
%  FITNESS FOR A PARTICULAR PURPOSE.
%
%  The main file must define before \input-ing this file:
%    \booktitle, \bookshort, \bookauthor, \bookid (ml or dl)
% =====================================================================

%%fakesection Load packages (from Napkin / evan.sty)
\usepackage{lmodern}
\usepackage[T1]{fontenc}
\usepackage[utf8]{inputenc}
\usepackage[pdfusetitle]{hyperref}
\usepackage{amsmath,amssymb,amsthm}
\usepackage{mathrsfs}
\usepackage[usenames,svgnames,dvipsnames,table]{xcolor}
\usepackage{textcomp}
\usepackage{enumerate}
\usepackage{mathtools}
\usepackage{microtype}
\usepackage[normalem]{ulem}
\usepackage{stmaryrd}
\usepackage{multirow}
\usepackage{booktabs}
\usepackage{array}
\usepackage{tabularx}
\usepackage{bm}
\usepackage{graphicx}
\usepackage[nameinlink]{cleveref}
\usepackage{etoolbox}

%%fakesection Napkin pedagogy commands
\newcommand{\prototype}[1]{%
	\emph{{\color{red} Prototypical example for this section:} #1} \par\medskip
}
\newcommand{\vocab}[1]{\textbf{\color{blue} #1}}
% \stub{...}: grey placeholder text while drafting.  None may survive
% into a final PDF; the read-through greps for the word STUB.
\newcommand{\stub}[1]{{\color{gray}\sffamily [STUB: #1]}}
\newcommand{\ii}{\item}
\newcommand{\half}{\frac{1}{2}}
\newcommand{\eps}{\varepsilon}
\newcommand{\defeq}{\coloneqq}
\newcommand{\inv}{^{-1}}
\newcommand{\wh}{\widehat}
\newcommand{\wt}{\widetilde}
\newcommand{\ol}{\overline}
\newcommand{\abs}[1]{\left\lvert #1 \right\rvert}
\newcommand{\norm}[1]{\left\lVert #1 \right\rVert}
\newcommand{\floor}[1]{\left\lfloor #1 \right\rfloor}
\newcommand{\ceiling}[1]{\left\lceil #1 \right\rceil}
\newcommand{\CC}{\mathbb C}
\newcommand{\NN}{\mathbb N}
\newcommand{\RR}{\mathbb R}
\newcommand{\ZZ}{\mathbb Z}
\newcommand{\EE}{\mathbb E}
\DeclareMathOperator*{\argmin}{arg\,min}
\DeclareMathOperator*{\argmax}{arg\,max}
\DeclareMathOperator{\rank}{rank}
\DeclareMathOperator{\Tr}{Tr}
\DeclareMathOperator{\sign}{sign}

%%fakesection Links
\hypersetup{
	linkcolor={red!50!black},
	citecolor={green!50!black},
	urlcolor={blue!80!black},
	pdfkeywords={machine learning, deep learning, placements, interviews},
	pdfauthor={\bookauthor},
	pdftitle={\booktitle},
	colorlinks,
}

%%fakesection TikZ and pgfplots (all figures are drawn, never pasted)
\usepackage{tikz}
\usetikzlibrary{arrows,arrows.meta,positioning,calc,shapes,shapes.geometric,
	fit,backgrounds,decorations.pathreplacing,matrix,patterns,shadows}
\usepackage{pgfplots}
\pgfplotsset{compat=1.18}
\usepgfplotslibrary{fillbetween,groupplots,colormaps}
\pgfplotsset{
	every axis/.append style={font=\small, line width=0.7pt,
		tick style={line width=0.5pt}},
	napkinplot/.style={width=0.82\linewidth, height=0.5\linewidth,
		grid=major, grid style={gray!20}, axis lines=left,
		legend style={draw=none, fill=white, fill opacity=0.8,
			text opacity=1, font=\footnotesize},
		cycle list name=napkincolors},
}
\pgfplotscreateplotcyclelist{napkincolors}{%
	{blue!80!black, mark=none, thick},
	{red!80!black, mark=none, thick, dashed},
	{ForestGreen, mark=none, thick, densely dotted},
	{orange!90!black, mark=none, thick, dashdotted},
	{violet, mark=none, thick},
	{gray, mark=none, thick, dashed}}

%% overfull boxes up to 2pt are tolerated (done criterion), wider ones are reported
\hfuzz=2pt
\vfuzz=2pt

%%fakesection Page layout (Napkin)
\usepackage[headsepline]{scrlayer-scrpage}
\renewcommand{\headfont}{}
\addtolength{\textheight}{3.14cm}
\setlength{\footskip}{0.5in}
\setlength{\headsep}{10pt}
\automark[chapter]{chapter}
\rohead{\footnotesize\thepage}
\rehead{\footnotesize \textbf{\sffamily \bookshort}, by \emph{\bookauthor}}
\lehead{\footnotesize\thepage}
\lohead{\footnotesize \leftmark}
\chead{}
\rofoot{}
\refoot{}
\lefoot{}
\lofoot{}

%%fakesection Fancy section and chapter heads (Napkin)
\renewcommand*{\sectionformat}{\color{purple}\S\thesection\autodot\enskip}
\renewcommand*{\subsectionformat}{\color{purple}\S\thesubsection\autodot\enskip}
\newcommand{\problemhead}{A few problems to think about}
\renewcommand{\thesubsection}{\thesection.\roman{subsection}}
\addtokomafont{chapterprefix}{\raggedleft}
\RedeclareSectionCommand[beforeskip=0.5em]{chapter}
\renewcommand*{\chapterformat}{%
	\mbox{\scalebox{1.5}{\chapappifchapterprefix{\nobreakspace}}%
		\scalebox{2.718}{\color{purple}\thechapter\autodot}\enskip}}
\addtokomafont{partprefix}{\rmfamily}
\renewcommand*{\partformat}{\color{purple}\scalebox{2.5}{\thepart}}

%%fakesection Theorems (Napkin tcolorbox + thmtools setup)
\usepackage{tcolorbox}
\tcbuselibrary{breakable,skins,hooks}
% patch tcolorbox to support continuation of a paragraph after a box
% (as in Napkin; from tex.stackexchange.com/questions/568782)
\makeatletter
\tcbset{
	after app={%
			\ifx\tcb@drawcolorbox\tcb@drawcolorbox@breakable
			\else
				\@endparenv
			\fi
		}
}
\appto\tcb@use@after@lastbox{\@endparenv\@doendpe}
\makeatother

\usepackage{thmtools}
\theoremstyle{definition}
\declaretheoremstyle[
	headfont=\sffamily\bfseries\color{MidnightBlue},
	headpunct={\\[3pt]},
	postheadspace={0pt},
]{thmtheorem}
\declaretheoremstyle[
	headfont=\bfseries\color{RawSienna},
	headpunct={\\[3pt]},
	postheadspace={0pt},
]{thmexample}

\tcbset{
	theorem box/.style={
			enhanced, arc=9pt, outer arc=10pt, colframe=blue,
			colback=TealBlue!5, boxrule=1pt, before skip=12pt,
			after skip=14pt, left=10pt, right=10pt, breakable
		},
	remark box/.style={
			boxrule=0pt, frame hidden, sharp corners, enhanced,
			borderline west={2pt}{0pt}{ForestGreen}, before skip=8pt,
			colback=ForestGreen!5, after skip=12pt, breakable,
			left=10pt, right=10pt
		},
	example box/.style={
			enhanced, sharp corners, arc=9pt, outer arc=10pt,
			colframe=RawSienna, colback=Salmon!5, boxrule=0.5pt,
			before skip=12pt, after skip=14pt, breakable, top=6pt,
			bottom=8pt, left=10pt, right=10pt
		},
	ques box/.style={
			boxrule=0pt, frame hidden, enhanced, sharp corners,
			before skip=8pt, after skip=12pt,
			borderline west={3pt}{0pt}{black},
			colback=RedViolet!5!gray!5, breakable, left=10pt, right=10pt
		}
}

\declaretheorem[style=thmtheorem,name=Theorem,numberwithin=section]{theorem}
\tcolorboxenvironment{theorem}{theorem box}
\declaretheorem[style=thmtheorem,name=Lemma,sibling=theorem]{lemma}
\tcolorboxenvironment{lemma}{theorem box}
\declaretheorem[style=thmtheorem,name=Proposition,sibling=theorem]{proposition}
\tcolorboxenvironment{proposition}{theorem box}
\declaretheorem[style=thmtheorem,name=Corollary,sibling=theorem]{corollary}
\tcolorboxenvironment{corollary}{theorem box}
\declaretheorem[style=thmexample,name=Example,sibling=theorem]{example}
\tcolorboxenvironment{example}{example box}

\declaretheoremstyle[
	headfont=\bfseries\sffamily\color{ForestGreen!70!black},
	bodyfont=\normalfont,
	headpunct={ --- },
]{thmremark}
\declaretheoremstyle[
	headfont=\bfseries,
	bodyfont=\normalfont\small
]{thmques}
\declaretheorem[name=Question,sibling=theorem,style=thmques]{ques}
\tcolorboxenvironment{ques}{ques box}
\declaretheorem[name=Exercise,sibling=theorem,style=thmques]{exercise}
\tcolorboxenvironment{exercise}{ques box}
\declaretheorem[name=Remark,sibling=theorem,style=thmremark]{remark}
\tcolorboxenvironment{remark}{remark box}

\theoremstyle{definition}
\newtheorem{claim}[theorem]{Claim}
\newtheorem{definition}[theorem]{Definition}
\newtheorem{fact}[theorem]{Fact}
\newtheorem{abuse}[theorem]{Abuse of Notation}

% Moral (Napkin): one-line takeaway in a green frame
\newenvironment{moral}{%
	\begin{tcolorbox}[boxrule=0.4pt,colframe=green!70!black,sharp corners,
				standard jigsaw,opacityback=0,left=10pt,right=10pt,top=3pt,bottom=3pt,
				before skip=10pt,after skip=20pt]%
			\bfseries\color{green!50!black}}%
		{\end{tcolorbox}}

%%fakesection Problems with hints and solutions (Napkin, via answers.sty)
\declaretheorem[style=definition,name=Problem,numberwithin=chapter,
	postheadhook={\phantomsection\label{prob:\theproblem}}]{problem}
\renewcommand{\theproblem}{\thechapter\Alph{problem}}
\usepackage{answers}
\Newassociation{hint}{hintitem}{\bookid-hints}
\Newassociation{sol}{solitem}{\bookid-solns}
\renewcommand{\solutionextension}{out}
\renewenvironment{hintitem}[1]{\item[{\bfseries\hyperref[prob:#1]{#1}.}]%
	\phantomsection\label{hint:#1}}{}
\renewenvironment{solitem}[1]{\item[{\bfseries\hyperref[prob:#1]{#1}.}]%
	\phantomsection\label{sol:#1}}{}
% Every problem gets an anchor and, at its end, links to hint and solution.
\AtEndEnvironment{problem}{\par\hfill{\footnotesize\sffamily
		\hyperref[hint:\theproblem]{hint}\,$\cdot$\,%
		\hyperref[sol:\theproblem]{solution}}}
% Problems in the big end-of-book set are numbered P1, P2, ...
\newcommand{\bigproblemset}{\renewcommand{\theproblem}{P\arabic{problem}}}

%%fakesection Chili difficulty icons (drawn in TikZ instead of Napkin's PNG)
\newcommand{\chili}{%
	\begin{tikzpicture}[baseline=-0.6ex, x=1em, y=1em]
		\fill[red!85!black] (0,0) .. controls (0.35,0.35) and (0.9,0.25) ..
		(1.35,-0.35) .. controls (1.0,-0.05) and (0.45,-0.05) .. (0.1,-0.25)
		.. controls (0.0,-0.18) .. cycle;
		\fill[white, opacity=0.5] (0.25,0.06) .. controls (0.55,0.2) ..
		(0.9,0.08) .. controls (0.55,0.12) .. cycle;
		\draw[ForestGreen!80!black, line width=1.3pt, line cap=round]
		(0.02,-0.08) -- (-0.18,0.05) .. controls (-0.28,0.12) .. (-0.25,0.3);
	\end{tikzpicture}}
\newcommand{\prechili}{\vspace*{0.3em}\hspace*{\fill}}
\newcommand{\onechili}{\marginpar{\prechili\chili}}
\newcommand{\twochili}{\marginpar{\prechili\chili\chili}}
\newcommand{\threechili}{\marginpar{\prechili\chili\chili\chili}}

%%fakesection Table of contents (Napkin)
\makeatletter
\pretocmd{\tableofcontents}{%
	\if@openright\cleardoublepage\else\clearpage\fi
	\pdfbookmark[0]{\contentsname}{toc}%
}{}{}%
\makeatother
\setcounter{tocdepth}{1}
\RedeclareSectionCommand[tocnumwidth=4.2em]{part}
\RedeclareSectionCommand[tocpagenumberwidth=2.8em,tocnumwidth=4.2em]{chapter}
\RedeclareSectionCommand[tocpagenumberwidth=2.8em,tocnumwidth=2.8em]{section}

%%fakesection Bibliography (Napkin style)
\usepackage[backend=biber,backref=true,style=alphabetic,maxbibnames=6,
	maxcitenames=2]{biblatex}
\DeclareLabelalphaTemplate{
	\labelelement{
		\field[final]{shorthand}
		\field{label}
		\field[strwidth=2,strside=left]{labelname}
	}
	\labelelement{
		\field[strwidth=2,strside=right]{year}
	}
}
\DeclareFieldFormat{labelalpha}{\textbf{\scriptsize #1}}
\DefineBibliographyStrings{english}{%
	backrefpage  = {cited p.},
	backrefpages = {cited pp.}
}
\DeclareFieldFormat{journaltitle}{\mkbibemph{#1},}
\DeclareFieldFormat[inbook,thesis]{title}{\mkbibemph{#1}\addperiod}
\DeclareFieldFormat[article]{title}{#1}
\DeclareFieldFormat[article]{volume}{\textbf{#1}\addcolon\space}
\renewcommand{\mkbibnamegiven}[1]{\textsc{#1}}
\renewcommand{\mkbibnamefamily}[1]{\textsc{#1}}
\renewcommand{\mkbibnameprefix}[1]{\textsc{#1}}
\renewcommand{\mkbibnamesuffix}[1]{\textsc{#1}}
\renewcommand{\finentrypunct}{}

%%fakesection Mini ToC (Napkin)
\usepackage[tight]{minitoc}
\mtcsetfont{parttoc}{chapter}{\sffamily\bfseries}
\mtcsetfont{parttoc}{section}{\footnotesize\rmfamily\upshape\mdseries}
\mtcsetfont{parttoc}{subsection}{\footnotesize\rmfamily\upshape\mdseries}
\setcounter{parttocdepth}{1}
\renewcommand*{\partheadstartvskip}{\vspace*{20em}}
\renewcommand*{\partheadendvskip}{}
\renewcommand\beforeparttoc{\noindent{\bfseries \Large Part \thepart: Contents}}
\makeatletter\renewcommand\@pnumwidth{2.4em}\makeatother
\doparttoc[n]

% =====================================================================
% =====================================================================
%%  ML / DL MACROS  (new code for these books, not from Napkin)
% =====================================================================
% =====================================================================

%% ---- generated values: every number in a worked example comes from
%% ---- code/ via build.sh, which writes \gvset{script/key}{value} lines
%% ---- into gen/<book>/*.tex.  \gv{script/key} prints the value; an
%% ---- undefined key prints a red ?? and a LaTeX warning.
\makeatletter
\newcommand{\gvset}[2]{\expandafter\gdef\csname gv@#1\endcsname{#2}}
\newcommand{\gv}[1]{%
	\ifcsname gv@#1\endcsname\csname gv@#1\endcsname
	\else{\color{red}\bfseries ??[#1]}%
		\PackageWarning{gv}{Undefined generated value `#1'}\fi}
\makeatother

%% ---- code listings and real program output
\usepackage{listings}
\usepackage{fancyvrb}
\definecolor{codebg}{RGB}{248,248,245}
\definecolor{codekw}{RGB}{0,70,160}
\definecolor{codecm}{RGB}{110,110,110}
\definecolor{codestr}{RGB}{160,60,0}
\lstdefinestyle{napkinpy}{
	language=Python, basicstyle=\ttfamily\footnotesize,
	keywordstyle=\color{codekw}\bfseries, commentstyle=\color{codecm}\itshape,
	stringstyle=\color{codestr}, showstringspaces=false,
	columns=fullflexible, keepspaces=true, breaklines=true,
	breakatwhitespace=false, postbreak=\mbox{\textcolor{gray}{$\hookrightarrow$}\space},
	numbers=left, numberstyle=\tiny\color{gray}, numbersep=6pt,
	xleftmargin=14pt, frame=none, upquote=true, tabsize=4,
	literate={≈}{{$\approx$}}1 {→}{{$\to$}}1 {—}{{---}}1 {λ}{{$\lambda$}}1
	{σ}{{$\sigma$}}1 {μ}{{$\mu$}}1 {θ}{{$\theta$}}1 {η}{{$\eta$}}1
	{ε}{{$\varepsilon$}}1 {Σ}{{$\Sigma$}}1 {α}{{$\alpha$}}1 {β}{{$\beta$}}1
	{γ}{{$\gamma$}}1 {≥}{{$\geq$}}1 {≤}{{$\leq$}}1 {×}{{$\times$}}1
	{²}{{$^2$}}1 {·}{{$\cdot$}}1 {∂}{{$\partial$}}1,
}
\lstset{style=napkinpy}
\tcbuselibrary{listings}
\newtcbinputlisting{\codefile}[2][]{%
	listing file={#2}, listing only, enhanced, breakable,
	colback=codebg, colframe=gray!40, boxrule=0.4pt, sharp corners,
	left=2pt, right=2pt, top=1pt, bottom=1pt,
	listing options={style=napkinpy}, title={\ttfamily\footnotesize #1},
	coltitle=black, colbacktitle=gray!15, fonttitle=\footnotesize,
	before skip=8pt, after skip=10pt}
\newtcblisting{pycode}[1][]{%
	listing only, enhanced, breakable, colback=codebg, colframe=gray!40,
	boxrule=0.4pt, sharp corners, left=2pt, right=2pt, top=1pt, bottom=1pt,
	listing options={style=napkinpy}, before skip=8pt, after skip=10pt, #1}
% \outputfile{path}: verbatim output captured from a script run.
\newcommand{\outputfile}[1]{%
	\begin{tcolorbox}[enhanced, breakable, colback=black!80, colupper=white,
			boxrule=0pt, sharp corners, left=4pt, right=4pt, top=2pt, bottom=2pt,
			before skip=6pt, after skip=10pt,
			title={\ttfamily\scriptsize output (generated by build.sh)},
			fonttitle=\scriptsize, colbacktitle=black!60, coltitle=white]
		\VerbatimInput[fontsize=\footnotesize, formatcom=\color{white}]{#1}
	\end{tcolorbox}}

%% ---- small algorithm box
\newtcolorbox{algo}[1][]{enhanced, breakable, colback=white,
	colframe=MidnightBlue!60, boxrule=0.6pt, sharp corners,
	fonttitle=\sffamily\bfseries\small, coltitle=white,
	colbacktitle=MidnightBlue!70, title={Algorithm: #1},
	before skip=10pt, after skip=12pt, left=6pt, right=6pt}

%% ---- "Trap" marker used at the start of remarks about wrong MCQ answers
\newcommand{\trap}{\textbf{\color{red!70!black}\sffamily Trap.}\ }
\newcommand{\oatag}{\textsuperscript{\color{orange!80!black}\sffamily\scriptsize OA}}

%% ---- notation
\renewcommand{\vec}[1]{\bm{#1}}
\newcommand{\x}{{\bm{x}}}
\newcommand{\y}{{\bm{y}}}
\newcommand{\w}{{\bm{w}}}
\newcommand{\bb}{{\bm{b}}}
\newcommand{\z}{{\bm{z}}}
\newcommand{\h}{{\bm{h}}}
\newcommand{\uu}{{\bm{u}}}
\newcommand{\vv}{{\bm{v}}}
\newcommand{\bmu}{{\bm{\mu}}}
\newcommand{\btheta}{{\bm{\theta}}}
\newcommand{\bbeta}{{\bm{\beta}}}
\newcommand{\balpha}{{\bm{\alpha}}}
\newcommand{\bSigma}{{\bm{\Sigma}}}
\newcommand{\X}{{\mathbf{X}}}
\newcommand{\Y}{{\mathbf{Y}}}
\newcommand{\W}{{\mathbf{W}}}
\newcommand{\I}{{\mathbf{I}}}
\newcommand{\A}{{\mathbf{A}}}
\newcommand{\B}{{\mathbf{B}}}
\newcommand{\Hm}{{\mathbf{H}}}
\newcommand{\Q}{{\mathbf{Q}}}
\newcommand{\K}{{\mathbf{K}}}
\newcommand{\V}{{\mathbf{V}}}
\newcommand{\U}{{\mathbf{U}}}
\newcommand{\M}{{\mathbf{M}}}
\newcommand{\zero}{{\bm{0}}}
\newcommand{\one}{{\bm{1}}}
\newcommand{\T}{^{\top}}
\newcommand{\E}{\mathbb{E}}
\newcommand{\Prob}{\mathbb{P}}
\newcommand{\Normal}{\mathcal{N}}
\newcommand{\Loss}{\mathcal{L}}
\newcommand{\D}{\mathcal{D}}
\newcommand{\Hyp}{\mathcal{H}}
\newcommand{\ind}[1]{\mathbb{1}\!\left[#1\right]}
\newcommand{\given}{\,\middle|\,}
\newcommand{\dd}{\,\mathrm{d}}
\newcommand{\pd}[2]{\frac{\partial #1}{\partial #2}}
\newcommand{\iid}{\overset{\text{i.i.d.}}{\sim}}
\DeclareMathOperator{\Var}{Var}
\DeclareMathOperator{\Cov}{Cov}
\DeclareMathOperator{\Corr}{Corr}
\DeclareMathOperator{\Bias}{Bias}
\DeclareMathOperator{\MSE}{MSE}
\DeclareMathOperator{\KL}{KL}
\DeclareMathOperator{\diag}{diag}
\DeclareMathOperator{\softmax}{softmax}
\DeclareMathOperator{\relu}{ReLU}
\DeclareMathOperator{\logit}{logit}
\DeclareMathOperator{\Ber}{Bernoulli}
\DeclareMathOperator{\Bin}{Bin}
\DeclareMathOperator{\Poi}{Poisson}
\DeclareMathOperator{\Unif}{Unif}
\DeclareMathOperator{\vect}{vec}
\DeclareMathOperator{\conv}{conv}
\newcommand{\sigmoid}{\sigma}

%% ---- references in plain text to the companion book
\newcommand{\otherbook}[1]{\emph{#1}}

%% ---- backmatter lists of hints/solutions (robust when still empty)
\newcommand{\printanswers}[1]{%
	\begin{enumerate}\setlength{\itemsep}{3pt}%
		\item[]\vspace{-1.5\baselineskip}%
		\input{#1}%
	\end{enumerate}}

%% ---- file paths in running text (breakable, underscores allowed)
\newcommand{\script}[1]{\nolinkurl{code/\bookid/#1.py}}
\newcommand{\fpath}[1]{\nolinkurl{#1}}
% \code{max_depth}: inline code that may break at _ . / - (no escaping needed)
\makeatletter\g@addto@macro\UrlBreaks{\do\_\do\.\do\-\do\=}\makeatother
\newcommand{\code}[1]{\nolinkurl{#1}}
